% =============================================================================
% 2.8 Stratégie d'évaluation et de validation
% =============================================================================
\section{Stratégie d'évaluation et de validation}
\label{sec:ch2-evaluation}

\subsection{Évaluation de la récupération}

L'évaluation du moteur de recherche s'appuie sur un jeu de référence
(\anglais{golden set}) de trente questions, dont les extraits pertinents
auront été identifiés manuellement par des praticiens. Les métriques
retenues sont le rappel à 20, la précision à 5 et le rang réciproque
moyen.

Quatre configurations seront comparées~: recherche dense seule, BM25 seul,
hybride par fusion de rangs, hybride suivi d'un ré-ordonnancement. Les
seuils d'interprétation sont les suivants~: un rappel à 20 inférieur à
$0{,}70$ signale un problème d'ingestion ou de découpage~; un rappel
supérieur à $0{,}80$ indique un système opérationnel.

\subsection{Évaluation de la génération}

L'évaluation de la génération est confiée à des experts du domaine, sur
un dossier de test. Chaque section est notée selon quatre critères~:
fidélité au contexte, ancrage des affirmations, exhaustivité au regard
des exigences, exactitude factuelle. Un protocole complémentaire, de type
modèle juge (RAGAS), accélère l'itération mais ne substitue pas le
jugement expert.

Les seuils d'alerte sont une fidélité inférieure à $0{,}80$ (problème de
contexte ou d'invite) et un ancrage inférieur à $0{,}70$ (problème de
citation).

\subsection{Validation des hypothèses}

Les hypothèses H1 à H5, formulées en introduction générale, se traduisent
en comparaisons contrôlées.

\begin{btptable}{Validation empirique des hypothèses}
\label{tab:validation-h}
\begin{tabularx}{\linewidth}{@{}C{1.1cm} Y L{3.8cm} C{1.8cm}@{}}
\toprule
\rowcolor{btppale}
\btpth{Réf.} & \btpth{Énoncé opérationnel} & \btpth{Comparaison} & \btpth{Gain visé} \\
\midrule
H1 & La recherche hybride surpasse la recherche dense seule. & Rappel à 20 & $+10$\,\% \\
H2 & Le découpage parent-enfant surpasse la fenêtre fixe. & Ancrage et exhaustivité & $+15$\,\% \\
H3 & Le ré-ordonnancement surpasse la seule fusion de rangs. & Précision à 5 & $+20$\,\% \\
H4 & La génération par section surpasse la génération globale. & Couverture des exigences & $+25$\,\% \\
H5 & La revue humaine améliore la fiabilité sans annuler le gain. & Taux de validation, temps de cycle & $+30$\,\% \\
\bottomrule
\end{tabularx}
\end{btptable}

\subsection{Robustesse}

Cinq dispositifs bornent le comportement du système en cas de défaillance.

\begin{enumerate}
  \item des validations déterministes avant et après chaque appel au
        modèle~;
  \item des reprises bornées (deux essais au plus) avant interruption
        humaine~;
  \item des délais d'expiration sur les travaux longs, avec arrêt propre~;
  \item des transactions pour la cohérence des écritures duales (relationnel
        et vectoriel)~;
  \item des journaux structurés, exploitables pour l'audit et le diagnostic.
\end{enumerate}
